%!TEX root = main.tex

%The proof of the functional correctness of ExpandConnects and ExpandWhens.

At first, we give the specifications of ExpandConnects and ExpandWhens, then prove that the implementations conform to the specifications. 


The correctness proof for ExpandWhens compares the Coq implementation with (a formalization of) the output specified above.
It requires a mutual induction over statements and statement sequences
because the \lstinline~when~ statement contains two statement sequences,
but the proof is more long than complicated.

\subsection{Correctness of ExpandConnects}
Recall the Expand Connects pass is aimed to decompose the connect statement of aggregate
types to ground type connections.  Let $ce^{in}$ be the components
environment, and $ss^{in}$ be the list of connect statements in the input
circuit diagram $D_M^{in}$.
Let $ss^{out}$ be the output list of connect statements after
applying the Expand Connections pass.
The input $D_M^{in}$ is required to be the result of Infer Types pass, in order to make sure the
connections are made between equivalent types and the components are already
recorded in the environment.

We state the correctness upon the types of the left-hand side and the right-hand side
expressions in the new generated connect statements.
First, there are only ground types in $ss'$,
%% \begin{lemma}
%% \forall ce, 
%% \end{lemma}


\xiaomu{Stopped}


\subsection{Correctness of ExpandWhens}
We have formalized the following central statement of the ExpandWhens transformation:

\begin{quote}
For every ground-type writable (sub)component,
ExpandWhens produces the connect or \prettyll{is invalid} statement
that appears last on any execution path through the module.
\linebreak
If there are multiple execution paths where the (sub)component is declared,
ExpandWhens seleects between them using a multiplexer or \prettyll{validif} expression.
\end{quote}

To describe the execution paths,
we introduce an auxiliary notion:
\begin{definition}
An \emph{expression tree} is a binary tree whose internal nodes are annotated with conditions
(usually, conditions found in \prettyll{when} statements)
and whose leaves are annotated with expressions or one of three special constants:
{\tt invalid}, {\tt not connected}, or {\tt undeclared}.
\end{definition}
An expression tree is able to express the declarations and connects over all execution paths of one (sub)component.
This tree serves as the basis to verify
whether the correct expression is actually stored for that variable in the connect map generated by {\tt expand\_branch}.

The proof is tedious because it requires a mutual recursion over statement sequences and statements,
because the \prettyll{when} statement again contains two statement sequences.
It also requires to regenerate some parts of the component environment to find components
that are only declared in a branch of some \prettyll{when c} statement;
connect statements for such components may only appear within that same branch,
and the condition \prettyll{c} is not checked in multiplexers for them.

%%%%% the formal semantics of HiFIRRTL removed %%%%%%%%%%
%%%%% the formal semantics of HiFIRRTL removed %%%%%%%%%%
\hide{
The formal semantics of HiFIRRTL will be defined by the axiomatic semantics of the compilation passes, plus the operational semantics of LoFIRRTL.
%InferTypes, InferWidths, InferResets, ExpandConnects, ExpandWhens, LowerTypes, and RemoveReset.  
%
More precisely, the axiomatic semantics of each compilation pass $\pi$ will be specified by a formula $\varphi_{\pi}(\cdots)$, and 
%Note that the sequence of statements in $ss$ will keep unchanged in the compilation passes.
%
%ce1 = φ_InferType(ss) /\
%ce2 = φ_InferWidth(ss, ce1) /\
%ce3 = φ_InferReset(ss, ce2) /\
%(ss1, ce4) = φ_ExpandConnect(ss, ce3) /\
%(ss2, cm) = φ_ExpandWhen(ss1, ce4) /\
%(ss3, ce5, cm1) = φ_LowerTypes(ss2, ce4, cm) /\
%(ss4, ce6, cm2) = φ_RemoveResets(ss3, ce5, cm1)
the semantics of $M$ is defined by the formula
\[
\begin{array}{l c l}
\varphi_M & \eqdef & \varphi_{ResolveKinds}(ss, ce_1)  \wedge  \varphi_{InferTypes}(ss, ce_1, ce_2)\ \wedge \\
&& \varphi_{InferWidths}(ss, ce_2, ce_3) \wedge  \varphi_{InferResets}(ss, ce_3, ce_4)\ \wedge \\
&& \varphi_{ExpandConnects}(ss, ce_4, ss_1, ce_5) \wedge \varphi_{ExpandWhens}(ss_1, ce_5, ss_2, cm)\ \wedge \\
&& \varphi_{LowerTypes}(ss_2, ce_5, cm, ss_3, ce_6, cm_1) \wedge  \varphi_{RemoveReset}(ss_3, ce_6, cm_1, ss_4, ce_7, cm_2).
\end{array}
\]
Note that each formula $\varphi_{\pi}$ actually specifies the graph of a function corresponding to $\pi$. For instance, the formula $\varphi_{InferWidths}(ss, ce_2, ce_3)$ specifies a function that takes $ss$ and $ce_2$ as inputs and produces the output $ce_3$.  The final output $ss_4$ is actually a sequence of statements already in LoFIRRTL form, whose operational semantics has been defined in Section~\ref{sec-lofirrtl-sem}.

In the sequel, we shall illustrate the definition of the semantics of compilation passes by describing $\varphi_{InferWidths}$ and $\varphi_{ExpandWhens}$. The semantics of the other compilation passes can be found in Appendix~\ref{sec-appen-sem-cpass}.  

%Let us fix $ss = (s_1, \cdots, s_n)$ in the sequel.

$\varphi_{InferWidths}(ss, ce, ce') \eqdef \varphi_{InferWidths}(ss, ce, ce')$

 
$\varphi_{ExpandWhens}$

last connect semantics

four states of components $\{\bot, \sim, -\}$ (undefined, defined, invalid), 

For $v \in \varset(M)$, we construct an expression $EWRhsExp_{ss, ce, v}(rhs)$ inductively as follows,  
and for $i \in [n]$, 
\begin{itemize}
\item If $ss = \varepsilon$, then $EWRhsExp_{ss, ce, v}(rhs) = rhs$, 
%
\item If $ss = s \cdot tl$, then
\begin{itemize}
\item case $s \equiv v' \mbox{ is invalid}$: if $v = v'$ and $rhs = \bot$ or $\sim$, then $EWRhsExp_{ss, ce, v}(rhs) = EWRhsExp_{tl, v, -}$, otherwise, $EWRhsExp_{ss, v, rhs} = EWRhsExp_{tl, v, rhs}$,
%
\item case $s \equiv \mbox{wire } v': t$: if $v = v'$ and $rhs = \bot$, then $EWRhsExp_{ss, v, rhs} = EWRhsExp_{tl, v, \sim}$, otherwise,  $EWRhsExp_{ss, v, rhs} = EWRhsExp_{tl, v, rhs}$,
%
\item case $s \equiv \mbox{reg } v': t, c$: if $v = v'$ and $rhs = \bot$, then $EWRhsExp_{ss, v, rhs} = EWRhsExp_{tl, v, \sim}$, otherwise, $EWRhsExp_{ss, v}  = EWRhsExp_{tl, v, rhs}$, 
%
\item case $s \equiv v'<= e$: if $v = v'$, then $EWRhsExp_{ss, v, rhs} = EWRhsExp_{tl, v, e}$, otherwise,  $EWRhsExp_{ss, v, rhs} = EWRhsExp_{tl, v, rhs}$,
%
\item case $s \equiv \mbox{when } c: sst \mbox{ else } ssf$: suppose $rhst = EWRhsExp_{sst, v, rhs}$ and $rhsf = EWRhsExp_{ssf, v, rhs}$, then
$EWRhsExp_{ss, v, rhs} = EWRhsExp_{tl, v, mux(c, rhst, rhsf)}$,
%
%\item case $s \equiv v' \mbox{ inst of } v''$: if $v = v'$ and $rhs = \bot$, then $EWRhsExp_{ss, v, rhs} = EWRhsExp_{tl, v, \sim}$, otherwise, $EWRhsExp_{ss, v, rhs} = EWRhsExp_{tl, v, rhs}$,
%
\item for the other statements,  $EWRhsExp_{ss, v, rhs} = EWRhsExp_{tl, v, rhs}$.
%\begin{itemize}
%\item if $rhst = \bot$ or $\sim$, then $EWRhsExp_{ss, v, rhs} = EWRhsExp_{tl, v, rhst}$,
%
%\item otherwise, if $rhsf = \bot$, then $EWRhsExp_{ss, v, rhs} = EWRhsExp_{sst, v, rhs}$,
%
%\item if $rhst \neq \bot$ and $rhsf \neq \bot$, then if $rhst = rhsf$, then $EWRhsExp_{ss, v, rhs}  = rhst$, otherwise, $EWRhsExp_{ss, v, rhs} = EWRhsExp_{tl, v, mux(c, rhst, rhsf)}$,
%
%\item 
%\end{itemize}
\end{itemize}
\end{itemize}

Finally, we transform $EWRhsExp_{ss, v, rhs}$ into $EWRhsExp'_{ss, v, rhs}$ as follows: 
\begin{itemize}
\item If $EWRhsExp_{ss, v, rhs} = mux(c, e_1, e_2)$, then 
\begin{itemize}
\item if both $e_1$ or $e_2$ contains occurrences of $\bot, \sim$, then $EWRhsExp'_{ss, v, rhs} = \bot$, 
%
\item if both 
\end{itemize} 
\end{itemize}

Define a recursive function $expandWhens$
}
%%%%% the formal semantics of HiFIRRTL removed %%%%%%%%%%
%%%%% the formal semantics of HiFIRRTL removed %%%%%%%%%%
